\documentclass[]{report}

\usepackage{import}

\import{../}{settings.tex}

\title{Mathématiques financières}

\author{Cours de , transcrit par SADR TAHOURI Luca}
\date{}

\begin{document}

    \maketitle
    \tableofcontents

    \chapter{Introduction}

    \section{Taux d'intérêts}

    \section{Intérêts et capitalisation}

    \begin{df}{}{}
        On appelle "valeur présente" ou "valeur actuelle nette" (V.A.N.) ($F_n$ € à $t=n$). C'est un vocabulaire non utilisée dans le cadre "contingent" avec $F_n(\omega)$ "contingent" sera connu uniquement à la date future $t=n$.
    \end{df}

    \begin{re}{}{}
        Détenir $F_n$ € à $t=n$ équivaut à détenir $\frac{F_n}{(1+r)^n}$ € à $t=0$
    \end{re}

    Le prix d'un produit financier (a) (Par exemple une obligation) qui donne droit à recevoir $F_n$€ à $t=n$. Hypothèse: Il existe un compte d'épargne rémunéré au taux période $r$.\\
    Le prix de (a) à $t=0$ est la valeur actuelle nette: $\frac{F_n}{(1+r)^n}$€
\end{document}